% !PRINTED 637
\weberpage{637}
%<<CONT
of a \gesperrt{Gauss} sum belongs to the cyclotomic numbers,
and we have
\weq{24}{\varrho^m = \Pi^{t}\, \mathfrak{p}^t_{t'}.}

But by (9) and (11), if $m$ in them is replaced by $m_1$,
\weq{25}{\varrho_n^{m_1} = \pm\, p^{\frac{1}{2}\,m_1(m_2-1)+1}\, \psi_1(r_1^n)\,\dots\,\psi_{m_1-2}(r_1^n),}
\weq{26}{\varrho^n \varrho_n^{-1} = p^{\frac{1}{2}(n-1)(m_2-1)}\, \psi_1(r_1)\,\psi_2(r_1)\,\dots\,\psi_{n-1}(r_1)}
and thus $\varrho^n \varrho_n^{-1}$, $\varrho_n^{m_1}$, and consequently also $\varrho_n^m$, are numbers
of the field $\Omega_m$, and by (25) the combination $\varrho^n \varrho_n^{-1}$ occurring
in (26) remains contained in $\Omega_m$ when $n$ is increased by a
multiple of $m_1$. Accordingly $\varrho^n \varrho_n^{-1}$ is still a number in $\Omega_m$
even when $n$ does not have precisely the restricted meaning it
has had hitherto, but denotes an arbitrary number not divisible
by $q$.

The purpose of this consideration is to make \gesperrt{Kummer}'s
theorem (16) applicable simultaneously to several different primes
which belong to different values of $m_1$, and the result is then
expressed in the following theorem:

3. \gesperrt{Let $\varphi_n$ be a system of conjugate functionals of
the field $\Omega_m$ which is related to no other
natural primes than such as are congruent to 1
with respect to the modulus $q$ (or, for $q = 2$,
with respect to the modulus 4); further, let $t$
run through those numbers of the sequence
$1, 2, \dots, m - 1$ that are not divisible by $q$, and
let $t'$ be determined from the congruence
$t\,t' \equiv 1 \pmod{m}$.}

\gesperrt{Then the functional of the field $\Omega_m$:}
\[ \text{a)}\qquad \Phi_n = \Pi^{t}\, \varphi^t_{nt'} = \varepsilon\,\vartheta_n^m \]
\gesperrt{is associated with the $m$th power of a cyclotomic
number $\vartheta_n$, which is indeed itself contained in a
higher field, but whose $m$th power belongs to the
field $\Omega_m$, and}
\[ \text{b)}\qquad \vartheta_n^{-1}\vartheta_1^n = \alpha \]
\gesperrt{is a number of the field $\Omega_m$.}

The proof is contained in the foregoing discussions and
follows from the formulae (25), (26), if one
%CONT>>
